\section{Sharp bounds with at most eight impulses}\label{sec:impulses}

For a nonzero conservative admissible word, write a pair $(b,a)$ for
the successive impulses
\[
 +1,\underbrace{0,\ldots,0}_{b},-1,\underbrace{0,\ldots,0}_{a},
 \qquad b,a\in\N.
\]
A word of $q$ pairs has $2q$ nonzero impulses and displayed period
$p=\sum_{j=0}^{q-1}(2+b_j+a_j)$. The gaps have no imposed upper bound.
Write $B=\{0,1,2\}$. A \emph{base launch} is seed 1 or 2 at any phase,
or seed zero at a safe phase.

\subsection{From two gaps to one return map}

At the boundary immediately before $+1$, one pair acts by
\begin{equation}\label{eq:pair-return-map}
 Q_{b,a}(n)=T^a\bigl(T^{b+1}(T(n)+1)-1\bigr).
\end{equation}
After $+1$, the $b$ zero impulses and the negative-impulse step give
$b+1$ ordinary updates before the subtraction. The last $a$ updates
give a total of $2+b+a$ physical steps.

For the pair $(b,a)=(1,4)$ and seed $1$, the impulse word is
\[
 (+1,0,-1,0,0,0,0).
\]
One lap reads as follows, with each arrow labeled by the impulse
applied \emph{after} its $T$ update:
\begin{equation}\label{eq:pair14-first-lap}
 1\xrightarrow{+1}3\xrightarrow{0}5\xrightarrow{-1}7
 \xrightarrow{0}11\xrightarrow{0}17\xrightarrow{0}26
 \xrightarrow{0}13.
\end{equation}
Thus $Q_{1,4}(1)=13$, but the peak within the lap is $26$.
The census records both the boundary state and the interior peak.

The following bounds apply from a safe visit to $B$, even when the
gaps are arbitrarily long and the response spans several laps.

\begin{theorem}[Uniform bounds for base launches]\label{thm:impulse-bounds}
For $1\le q\le4$, every base launch for every word of $q$ pairs has value
at most the following bound and eventually enters a cycle of augmented
period at most the displayed multiple of $p$:
\begin{center}
\begin{tabular}{rrrr}\toprule
Nonzero impulses & $q$ & Maximum value & Maximum augmented period\\\midrule
2&1&122&$7p$\\
4&2&195020&$7p$\\
6&3&92779118&$12p$\\
8&4&54285893576&$13p$\\\bottomrule
\end{tabular}
\end{center}
Every listed maximum is attained by a positive base seed. Every period
bound is attained with primitive forcing and is then also the least
scalar response period. The same upper bounds apply to the future of
any trajectory after a safe visit to $B$; no claim about its earlier
values or the existence of such a visit is implicit.
\end{theorem}

\subsection{An exhaustive reduction of arbitrarily large gaps}

At a positive-impulse boundary, starting from $n$, set $u=T(n)+1$.
After $b$ zeros the negative impulse gives $T(T^b(u))-1$; the following
$a$ zeros complete the pair. Call the pair \emph{escaping} from $n$
when its endpoint exceeds 2. List all such pairs by the following finite
procedure whenever the encountered ordinary tails reach $B$:
\begin{enumerate}
\item For each $b\ge0$ while $T^b(u)>2$, put $v=T(T^b(u))-1$.
\item List each $a\ge0$ while $T^a(v)>2$, recording that endpoint.
\end{enumerate}
Once either ordinary tail enters $B$, larger gaps of that type cannot
produce an escaping pair: $T(B)\subseteq B$, and $T(w)-1\in B$ for
$w\in\{1,2\}$. From a base input a pair can escape only from 1;
inputs 0 and 2 give $u=1$ and $u=2$ and return to $B$.
Moreover, after a pair's initial positive impulse, any intermediate
visit to $B$ forces its endpoint to lie in $B$: the remaining zero
updates preserve $B$, and the sole negative update sends a positive
base value back into $B$. Thus a pair with endpoint outside $B$ has
not silently returned to $B$ in its interior.

\paragraph{All fifteen one-pair escaping choices.}
The first level of the procedure can be written out completely.
Starting from $1$, the value just after $+1$ is $u=3$, and its
ordinary tail is
\[
 3\longmapsto5\longmapsto8\longmapsto4\longmapsto2
 \longmapsto1\longmapsto2\longmapsto\cdots.
\]
Consequently the possible values immediately before the negative
impulse, and the remaining gap choices that escape, are
\begin{center}
\begin{tabular}{rlllr}\toprule
$b$ & Before $-1$ & Just after $-1$ & Escaping $a$ & Count\\\midrule
0&3&4&$a=0$&1\\
1&5&7&$0\le a\le9$&10\\
2&8&3&$0\le a\le3$&4\\
3&4&1&none&0\\
$b\ge4$&1 or 2&0 or 1&none&0\\\bottomrule
\end{tabular}
\end{center}
For the longest row, the ordinary tail after the negative impulse is
\[
 7,11,17,26,13,20,10,5,8,4,2,1,2,\ldots.
\]
Exactly its first ten entries exceed two. For $b=2$, the relevant
tail is $3,5,8,4,2,1,\ldots$, with four such entries. This gives
$1+10+4=15$ escaping pairs, including $(1,4)$ from
\eqref{eq:pair14-first-lap}.

The rows also settle infinite sets of omitted parameters. For example,
every pair $(1,a)$ with $a\ge10$ has endpoint
\[
 Q_{1,a}(1)=
 \begin{cases}2,&a\text{ even},\\1,&a\text{ odd}.\end{cases}
\]
There is no upper cutoff on $a$ in this assertion: reaching the exact
ordinary cycle $1\leftrightarrow2$ proves the formula for all remaining
gap lengths. The same principle is applied at every subsequent level
of the escaping-prefix construction.

To justify termination without assuming the ordinary Collatz conjecture,
form signed layers. Start with $B$, apply the next sign $+1$ or $-1$
after one $T$ step, and take its complete ordinary forward closure.
Explicitly, with $\mathcal L_0=B$,
\begin{equation}\label{eq:signed-layer-recurrence}
 \mathcal L_j=
 \bigcup_{x\in\mathcal L_{j-1}}
 \{T^a(T(x)+(-1)^{j+1}):a\ge0\}.
\end{equation}
The word ``signed'' refers to the alternating impulse signs. The
reachable values in these layers remain nonnegative. The first two
layers are small enough to display:
\begin{align*}
 \mathcal L_1&=\{1,2,3,4,5,8\},\\
 \mathcal L_2&=\{0,1,2,3,4,5,7,8,10,11,13,17,20,26\}.
\end{align*}
For instance, $T(B)+1=\{1,2,3\}$, and following its ordinary
tails gives $\mathcal L_1$. Applying $T(x)-1$ to this set produces
$\{0,1,3,4,7\}$; their ordinary tails give $\mathcal L_2$.
Every tail in these finitely many layers is explicitly followed until
it reaches $B$ or a previously verified tail; an unresolved tail or
integer overflow is a failure. The successive sizes and maxima are
\begin{center}\small
\begin{tabular}{rrrrrrrrr}\toprule
Signs applied&1&2&3&4&5&6&7&8\\\midrule
Size&6&14&85&317&1218&3279&10193&27870\\
Maximum&8&26&4616&15578&125252&642842&18534824&90272312\\\bottomrule
\end{tabular}
\end{center}
Each layer allows an arbitrary number of subsequent zero impulses.
Consequently these finite computations certify all ordinary tails
needed to enumerate up to four successive escaping pairs, for every
choice of gaps. They also bound the peaks during the corresponding
finite excursions. The numbers of rooted escaping prefixes are
\[
 15,\qquad 2854,\qquad 2416542,\qquad 2383451629.
\]
Here a rooted prefix starts from 1 immediately before its first positive
impulse and has endpoint greater than 2 after every complete pair.
These counts include rotations separately when their rootings qualify.
They are not counts of distinct cyclic forcing words.

Fix now any word of $q$ pairs and any base trajectory. Before the first
positive boundary the orbit remains in $B$. If it ever makes $q$
successive escaping pairs after departing $B$, that departure starts
from 1. The next $q$ pairs specify a rotation of the entire fixed word
and occur in the finite rooted-prefix list. Its entire future from
that departure is therefore the rooted fixed-word trajectory checked
in the core census. Before the first such excursion, each
departure lasts fewer than $q$ pairs, so the signed layers bound all
earlier peaks.

If there is no such excursion, every departure from $B$ returns within
$q$ pairs. At a fixed positive boundary, any state on the trajectory
is either in $B$ or is $r$ pairs after its most recent departure,
where $1\le r\le q-1$. For each $r$ the intervening pairs are fixed and
the departure seed must be 1, so there is at most one such outside
value. There are at most $3+(q-1)=q+2$ possibilities at that boundary.
Sampling once per word yields a repeated state and a cycle of period
at most $(q+2)p$. The peak is at most the largest signed-layer maximum
through $2q$ impulses. Both bounds are below the corresponding entries
of Theorem~\ref{thm:impulse-bounds}. This argument covers the entire
infinite complementary family, including gaps longer than any listed
escaping tail.

For example, take $q=2$ and suppose no departure from $1$ has two
consecutive escaping pairs. At a fixed positive boundary the orbit lies
in $B$ or is one pair past its last departure. That preceding pair is
fixed, so it contributes at most one more value, $Q_{b,a}(1)$. At most
four boundary values occur, giving an eventual period of at most four
laps for any gap lengths.

A signed layer bounds an excursion through a specified number of impulses.
A core word can run for several laps between base visits. The one-pair
example below has $\max\mathcal L_2=26$ and recurrent peak $122$.

\subsection{Exact core censuses}

For each rooted $q$-pair prefix, fix that same word periodically, rather
than allowing a new choice of gaps on later laps. Starting from its
root seed 1, compute its successive values at a fixed boundary and
record every intervening peak. The first repeated boundary value
identifies the eventual cycle and its period in displayed laps.
Any failure to close within the stated resource limit aborts the run;
the resource limit is never interpreted as proof of capture. For
$q\le3$, the included implementations additionally follow every seed
in $B$ at every positive boundary, allowing walks to join already
closed futures. This is redundant for the reduction but supplies an
extra check on phase coverage.

The one- and two-pair census is
\path{certificates/impulses/small.py}. The six-impulse census in
\path{certificates/impulses/independent_six/verify.cpp} closes all
2,416,542 rooted words, visiting a total of 32,037,568 distinct
positive-boundary states across the words and at most 44 for one word.
It checks unsigned integer overflow before forming $3n+1$ and skips
only the explicit ordinary cycles at 0 and $\{1,2\}$.

The eight-impulse core has 2,383,451,629 rooted words. It is partitioned
by the first escaping pair; the largest part is further split by a
third-prefix ordinal. Every partition must occur exactly once.
\path{certificates/impulses/full_excursions/census.cpp} closes every
word, using exact tables of finite ordinary iterates to reduce repeated
work. A second implementation,
\path{certificates/impulses/independent_eight/scan.cpp}, uses a different
partition and coast implementation. Both give maximum and recurrent
maximum 54,285,893,576 and maximum cycle length 13 laps. The stored
outputs and reproduction commands accompany the sources.

The quick checker validates partition sizes, histograms, excursion domains,
witness ownership and literal witness trajectories. Full reproduction
regenerates each partition and compares it with the stored certificate.
The upper bounds use the complete recurrence census together with the
signed layers and the proof for the infinite complement.

\subsection{A sharp witness in literal time}

Continue the fixed word $(1,4)$ from
\eqref{eq:pair14-first-lap}. The gaps are kept unchanged on every
lap. Its complete boundary cycle and the peaks of the corresponding
seven-step segments are
\begin{center}
\begin{tabular}{rrrr}\toprule
Lap & Input boundary value & Next boundary value & Segment peak\\\midrule
0&1&13&26\\
1&13&80&80\\
2&80&53&80\\
3&53&35&122\\
4&35&8&54\\
5&8&2&8\\
6&2&1&2\\\bottomrule
\end{tabular}
\end{center}
In the fourth row, the literal segment is
\[
 53,81,122,60,30,15,23,35.
\]
This supplies the peak $122$ without interpolation between sampled
values. The seven boundary inputs are distinct and the next is again
$1$, so this launch is already on an augmented cycle of $7\cdot7=49$
steps. The forcing has primitive period seven. Any scalar period
must be a multiple of that primitive forcing period, and a smaller
multiple would repeat a boundary value earlier. Hence the least
scalar response period is also $49$.

If the same seven-impulse list is displayed twice, the physical
recurrence and its scalar period stay unchanged. The displayed clock
then has length $14$, so a state must wait
$\operatorname{lcm}(49,14)=98$ steps to repeat both its scalar
value and that enlarged clock label. This explains the distinction
between a displayed augmented period and the scalar period attached
to primitive forcing. The sharpness witnesses in the theorem are
checked with their primitive clocks.

\begin{proof}[Proof of Theorem~\ref{thm:impulse-bounds}]
The signed-layer and complement arguments reduce every base launch to
the finite core or to smaller bounds. The complete core runs give the
displayed maxima and period bounds. For sharpness, the following
witnesses are replayed in physical time, including the forcing's
primitive period and the least scalar period of the eventual cycle.
\begin{center}\small
\begin{tabular}{rlll}\toprule
$q$&Quantity&Pairs in time order&Seed\\\midrule
1&peak and period&$(1,4)$&1\\
2&peak&$(1,2),(43,0)$&1\\
2&period&$(1,2),(10,2)$&1\\
3&peak&$(1,2),(23,11),(6,7)$&1\\
3&period&$(1,4),(2,1),(6,1)$&2\\
4&peak&$(1,2),(23,15),(14,32),(12,16)$&1\\
4&period&$(1,5),(1,0),(5,3),(0,2)$&1\\\bottomrule
\end{tabular}
\end{center}
Every displayed witness is launched immediately before the positive
impulse of its first listed pair. In general, if a scalar eventual cycle has
least period $r$, the recovered impulses
$d_k=n_{k+1}-T(n_k)$ have period $r$. For primitive forcing of period
$p$, this implies $p\mid r$, so the scalar cycle already returns to the
same forcing phase. Its least scalar and augmented periods agree.
\end{proof}

The census uses exact integer arithmetic. The tail exhaustion and
complement argument extend its bounds to arbitrary gaps for base launches.
